\documentclass[10pt,a4paper]{article}
\usepackage[T1]{fontenc}
\usepackage{lmodern,amsmath,amssymb,amsthm,mathtools,microtype,geometry,enumitem,hyperref}
\geometry{left=23mm,right=23mm,top=20mm,bottom=21mm}
\hypersetup{hidelinks,pdftitle={Foundations of Measure-Theoretic Probability - Part 1 - Revised Study Notes}}
\setlength{\parindent}{0pt}\setlength{\parskip}{4.5pt}
\setlist[itemize]{topsep=4pt,itemsep=3pt}
\newtheorem{proposition}{Proposition}
\newcommand{\N}{\mathbb N}\newcommand{\Q}{\mathbb Q}\newcommand{\R}{\mathbb R}
\begin{document}
\begin{center}
{\Large\bfseries PART 1: Measure-Theoretic Foundations}\par
{\normalsize Mathematical reading notes --- revised edition, 2026}
\end{center}
\vspace{8pt}
\textbf{Topics.} Motivations for the modern definition of probability; sigma-algebras and measurable sets (\S1.2); measures and the Dirac measure (\S1.3); monotonicity, continuity of measures and the union bound (\S1.4). The course's treatment of discrete probability mass functions (\S1.6) is not repeated here.

\section*{1. Why the classical descriptions of probability are not general definitions}
\subsection*{Classical probability}
For a finite set $S$ whose outcomes are \emph{equally likely}, the formula
\[
P(A)=\frac{|A|}{|S|}
\]
provides a useful model. However, it does not itself describe infinite sample spaces, and equal likelihood cannot simply be assumed for arbitrary experiments. In particular, a uniform probability distribution on \emph{all} positive integers does not exist: if every singleton had the same mass $c$, countable additivity would give $1=\sum_{n\ge1}c$, which is impossible for $c\ge0$.

\subsection*{Relative frequencies}
A familiar interpretation of probability is a limiting relative frequency,
\[
P(A)\approx \lim_{n\to\infty}\frac{\text{number of occurrences of }A\text{ in the first }n\text{ trials}}{n},
\]
when such a limit exists. This is an important consequence or interpretation of a probabilistic model, but it cannot replace the axioms in general: limits of observed frequencies may fail to exist, and functions obtained as limiting \emph{densities} on subsets of $\N$ need not be countably additive.

For example, the asymptotic density $d(A)=\lim_{n\to\infty}|A\cap\{1,\ldots,n\}|/n$, when it exists, gives $d(\{k\})=0$ for each singleton $\{k\}$, but $d(\N)=1$. If it were a probability measure on all subsets, countable additivity would imply $1=d(\N)=\sum_{k\ge1}d(\{k\})=0$. This is an example of the \emph{failure of countable additivity for asymptotic density}, not a genuine uniform random choice from $\N$.

\subsection*{Geometric probability and a nonmeasurable set}
For some finite intervals or regions, probability can be modeled by a ratio of lengths, areas or volumes. Such a model presupposes that the sets under discussion have a well-defined measure. In general one cannot consistently assign an extension of ordinary length to \emph{all} subsets of $\R$ while preserving translation invariance and countable additivity. The following Vitali-set construction illustrates why.

Define an equivalence relation on $[0,1]$ by
\[
x\sim y\quad\Longleftrightarrow\quad x-y\in\Q.
\]
The interval is partitioned into pairwise disjoint equivalence classes $[x]=\{y\in[0,1]:y-x\in\Q\}$. Each class is countable, because its elements form a subset of a rational translate; there are uncountably many classes, since a countable union of countable sets would be countable. By the axiom of choice, select one point from each class, and let $V\subset[0,1]$ be the set of representatives (a Vitali set).

Distinct points of $V$ cannot differ by a rational number. Consider \emph{rational} translates of $V$, indexed by $q\in\Q\cap[-1,1]$:
\[
V+q=\{v+q:v\in V\},\qquad
A=\bigcup_{q\in\Q\cap[-1,1]}(V+q).
\]
These translates are pairwise disjoint: if $v_1+q_1=v_2+q_2$, then $v_1-v_2=q_2-q_1\in\Q$, so by construction $v_1=v_2$ and hence $q_1=q_2$.

Every $x\in[0,1]$ differs by some rational $q\in[-1,1]$ from the representative of its class; therefore $x\in V+q$. Conversely $V\subset[0,1]$ and $q\in[-1,1]$, so each translate is contained in $[-1,2]$. Thus
\[
[0,1]\subseteq A\subseteq[-1,2].
\]
\begin{proposition}[Vitali contradiction]
There is no countably additive, translation-invariant, nonnegative measure $m$ defined on \emph{every subset} of $\R$ that assigns ordinary length to every interval.
\end{proposition}
\begin{proof}
If such a measure existed, monotonicity and the displayed inclusions would imply $1\le m(A)\le3$. On the other hand, using disjointness, translation invariance and countable additivity,
\[
m(A)=\sum_{q\in\Q\cap[-1,1]}m(V+q)
=\sum_{q\in\Q\cap[-1,1]}m(V).
\]
There are infinitely many such rationals. If $m(V)=0$ the sum is $0$; if $m(V)>0$ (including $+\infty$) the sum is $+\infty$. Both contradict $1\le m(A)\le3$.
\end{proof}
The argument \emph{does not} say that points of $A$ lying outside $[0,1]$ have measure zero, nor that $A$ has the same length as $[0,1]$. It motivates restricting the domain of measurable events instead of attempting to measure every subset.

\section*{2. Probability spaces and sigma-algebras (\S1.2)}
A \emph{sigma-algebra} $\mathcal S$ on a set $S$ is a collection of subsets satisfying
\begin{enumerate}[leftmargin=*,label=(\roman*)]
\item $\varnothing,S\in\mathcal S$;
\item $A\in\mathcal S\implies S\setminus A\in\mathcal S$;
\item $A_1,A_2,\ldots\in\mathcal S\implies \bigcup_{n\ge1}A_n\in\mathcal S$.
\end{enumerate}
By De Morgan's law it is also closed under countable intersections. Members of $\mathcal S$ are \emph{measurable sets}; $(S,\mathcal S)$ is a measurable space. Both the full power set $\mathcal P(S)$ and the trivial sigma-algebra $\{\varnothing,S\}$ are examples, for \emph{any} set $S$.

A \emph{probability space} is a triple $(S,\mathcal S,P)$, where $P:\mathcal S\to[0,1]$ is a probability measure. We do not ask $P$ to assign probabilities to all subsets of $S$ unless we have specifically chosen the full power set as the sigma-algebra.

\section*{3. Measures and probability measures (\S1.3)}
A measure on $(S,\mathcal S)$ is a map $\mu:\mathcal S\to[0,\infty]$ such that
\[
\mu(\varnothing)=0,
\qquad
\mu\left(\bigcup_{n\ge1}A_n\right)=\sum_{n\ge1}\mu(A_n)
\]
whenever the measurable sets $A_n$ are pairwise disjoint. The second property is \emph{countable additivity}. A finite measure satisfies $\mu(S)<\infty$; a probability measure satisfies $\mu(S)=1$. If $0<\mu(S)<\infty$, normalizing $\mu$ by $\mu(S)$ produces a probability measure.

\subsection*{The Dirac measure}
For a fixed point $x\in S$, define
\[
\delta_x(A)=\begin{cases}1,&x\in A,\\0,&x\notin A,\end{cases}
\qquad A\in\mathcal S.
\]
It is a probability measure: clearly $\delta_x(\varnothing)=0$ and $\delta_x(S)=1$. For pairwise disjoint $A_n$, either $x$ belongs to none of them, in which case both sides of countable additivity equal zero, or $x$ belongs to exactly one $A_k$, in which case both sides equal one.

\section*{4. Properties of measures (\S1.4)}
\subsection*{Monotonicity}
If $A\subseteq B$ are measurable, then $B=A\mathbin{\dot\cup}(B\setminus A)$. Consequently
\[
\mu(B)=\mu(A)+\mu(B\setminus A)\ge\mu(A).
\]
Thus measures are monotone.

\subsection*{Increasing and decreasing sequences of sets}
Write $A_n\uparrow A$ when $A_1\subseteq A_2\subseteq\cdots$ and $\bigcup_{n\ge1}A_n=A$. Write $A_n\downarrow A$ when $A_1\supseteq A_2\supseteq\cdots$ and $\bigcap_{n\ge1}A_n=A$. For example, for $n\ge1$,
\[
[0,1-1/n]\uparrow[0,1),\qquad
[-1/n,1]\downarrow[0,1].
\]
In the first example the right endpoint $1$ is \emph{not} part of the union; that matters for the set-theoretic statement even though the Lebesgue measures of $[0,1)$ and $[0,1]$ coincide.

\begin{proposition}[Continuity of measures]
For measurable sets, 
\begin{align*}
A_n\uparrow A&\implies\mu(A_n)\uparrow\mu(A),\\
A_n\downarrow A\ \text{and}\ \mu(A_1)<\infty
&\implies\mu(A_n)\downarrow\mu(A).
\end{align*}
\end{proposition}
\begin{proof}
For increasing sets, let $B_1=A_1$ and $B_n=A_n\setminus A_{n-1}$ for $n\ge2$. These sets are disjoint and satisfy $A_n=\dot\bigcup_{j=1}^n B_j$ and $A=\dot\bigcup_{j\ge1}B_j$. Countable additivity gives
\[
\lim_{n\to\infty}\mu(A_n)=\sum_{j\ge1}\mu(B_j)=\mu(A).
\]
For decreasing sets with $\mu(A_1)<\infty$, set $B_n=A_1\setminus A_n$ and $B=A_1\setminus A$. Then $B_n\uparrow B$. The first part gives $\mu(B_n)\uparrow\mu(B)$, and
\[
\mu(A_n)=\mu(A_1)-\mu(B_n)\longrightarrow
\mu(A_1)-\mu(B)=\mu(A).
\]
The finiteness assumption permits these subtractions.
\end{proof}
This result lets us compute the probability of limiting events by approximating them with simpler increasing or decreasing sequences of events.

\subsection*{The union bound}
For any finite or countable family of measurable sets $(A_i)$,
\[
\mu\left(\bigcup_{i\ge1} A_i\right)\le\sum_{i\ge1}\mu(A_i).
\]
For two sets $A,B$, since $A\cup B=A\mathbin{\dot\cup}(B\setminus A)$,
\[
\mu(A\cup B)=\mu(A)+\mu(B\setminus A)\le\mu(A)+\mu(B).
\]
Induction gives the finite version. For countably many sets, set $C_n=\bigcup_{i=1}^n A_i$, note $C_n\uparrow\bigcup_{i\ge1}A_i$, and use continuity from below:
\[
\mu\left(\bigcup_{i\ge1}A_i\right)
=\lim_{n\to\infty}\mu(C_n)
\le\lim_{n\to\infty}\sum_{i=1}^n\mu(A_i)
=\sum_{i\ge1}\mu(A_i).
\]
For disjoint sets equality holds by countable additivity; equality may also hold for overlapping sets when their overlaps have measure zero.

\section*{Reading note}
This part introduces the measure-theoretic language needed to define distributions and couplings in the subsequent parts. The Vitali construction is presented to explain why one works with a sigma-algebra of measurable events rather than arbitrary subsets.
\end{document}
